If $N$ is not a power of two, write $N=2^su$ with odd $u>1$. If $u\le 2^{s+1}+1$, use $u$ terms centered at $2^s$, beginning at $2^s-\dfrac{u-1}{2}\ge 0$. In the zero-endpoint equality case discard the zero; otherwise the start is positive. If $u>2^{s+1}+1$, use $2^{s+1}$ consecutive terms centered at $\dfrac{u}{2}$, beginning at $\dfrac{u-2^{s+1}+1}{2}>0$. Both constructions sum to $N$ and have at least two positive terms. Conversely, if $N$ is a sum of $k\ge 2$ consecutive positive integers beginning at $a$, then $2N=k(2a+k-1)$. The two factors have opposite parity. If $N$ were a power of two, the odd factor would have to be $1$. If $k$ is odd then $k\ge 3$; if $k$ is even then $2a+k-1\ge 3$. Both are impossible. This proves the characterization.
